% Pending complete proof, not yet integrated into checked chapter.
\begin{saresult}{A5.CVX.pending}{Proved under stated hypotheses}{The maximal locus for maps with convex fibres}
Fix an o-minimal expansion of an arbitrary real closed field $K$.
Let $X\subset K^n$, $Y\subset K^m$ be definable sets and let
$f:X\to Y$ be continuous, definable and definably proper. Assume every
nonempty fibre $F_y=f^{-1}(y)$ is a convex subset of $K^n$.
Then the maximal open definable-manifold-without-boundary locus over
which $f$ is definably locally trivial is definable and has nowhere
 dense complement.
\end{saresult}
\begin{proof}
A5.012 proves that definable properness makes $f$ a closed map on
closed definable subsets. In particular each $F_y$ is closed and bounded.
We use the dimension convention $\dim\varnothing=-1$. The function
$d:Y\to\{-1,0,\ldots,n\}$, $d(y)=\dim F_y$, is definable. Let $D\subset Y$
be its open definable locus of local constancy. Let $N\subset Y$ consist
of the points $y$ satisfying
\[
\begin{split}
 &\text{for all }x\in F_y\text{ and all }\delta>0\text{ there is }\varepsilon>0
 \text{ such that}\\
 &z\in Y,\quad \|z-y\|<\varepsilon
 \quad\Longrightarrow\quad
 \text{there is }x'\in F_z\text{ with }\|x'-x\|<\delta.
\end{split}                                                        \tag{A5.13.1}
\]
This is a first-order condition, so $N$ is definable. We claim that the
required maximal locus is
\[
                 U=Y_{\rm man}\cap D\cap\operatorname{Int}_Y(N).
                                                                    \tag{A5.13.2}
\]
Here $Y_{\rm man}$ is the open definable manifold locus from A5.005.
Any local product chart implies (A5.13.1) at every point of its inverse
image and makes $d$ constant. Thus every locally trivial manifold
neighbourhood is contained in the right-hand side.

For the converse, restrict to a definable open neighbourhood $V\subset Y$
of a point $y_0$ in the right side on which $d$ is constant and $V\subset N$.
The restriction $f_V:f^{-1}(V)\to V$ is definably proper. Formula (A5.13.1)
means exactly that $f_V$ is open: for an open subset $O\subset f^{-1}(V)$
and $x\in O$, choose a sufficiently small ball about $x$ contained in $O$
relative to the source, and apply (A5.13.1) at $f(x)$.
If $d=-1$, then $f^{-1}(V)=\varnothing$ and the empty product is a
trivialization. If $d=0$, a nonempty convex zero-dimensional set is a
singleton; the continuous bijective open map $f_V$ is a definable
homeomorphism, giving the claim. Hence assume $d\geq1$.

We first prove Hausdorff continuity of $y\mapsto F_y$ on $V$.
For nonempty closed bounded definable sets $A,B\subset K^n$, put
\[
 d_H(A,B)=\max\{\max_{a\in A}\operatorname{dist}(a,B),
                         \max_{b\in B}\operatorname{dist}(b,A)\}.
\]
All extrema exist by definable compactness, and the displayed expression
is definable in a definable family. Fix $y\in V$ and $\eta>0$.
The set $C=\{x\in f^{-1}(V):\operatorname{dist}(x,F_y)\geq\eta\}$
is closed in $f^{-1}(V)$, and $y\notin f_V(C)$. Closedness of the image
gives $F_z\subset\{x:\operatorname{dist}(x,F_y)<\eta\}$ for all $z$
sufficiently near $y$.
For the reverse inclusion, suppose it fails uniformly near $y$.
Definable curve selection and definable choice then provide continuous
definable curves $z(t)\to y$ and $a(t)\in F_y$ such that
$\operatorname{dist}(a(t),F_{z(t)})\geq\eta$ for $t>0$.
Definable compactness of $F_y$ gives an endpoint limit $a(t)\to a\in F_y$.
Openness at $a$ supplies, for all $z$ near $y$, a point of $F_z$ within
$\eta/3$ of $a$. For small $t$ it is within $2\eta/3$ of $a(t)$,
a contradiction. Thus $d_H(F_z,F_y)\to0$.

For each fixed $p\in K^n$ there is a unique nearest point
$P_y(p)\in F_y$. Existence follows from definable compactness. If two
minimizers were different, their midpoint would lie in $F_y$ by convexity
and have strictly smaller squared distance, by the parallelogram
identity. The map $y\mapsto P_y(p)$ is definable and continuous.
Indeed, along any definable curve $y(t)\to y$, Hausdorff continuity
bounds $P_{y(t)}(p)$ locally. Its endpoint limit $q$ belongs to $F_y$.
For every $a\in F_y$, Hausdorff continuity supplies points
$a(t)\in F_{y(t)}$ tending to $a$; the minimizing inequality passes
to the limit and gives $\|q-p\|\leq\|a-p\|$. Hence $q=P_y(p)$.
The definable curve criterion for continuity proves the assertion.

Choose affinely independent points $p_0,\ldots,p_d\in F_{y_0}$.
After shrinking $V$, the continuous points $p_i(y)=P_y(p_i)$ remain
affinely independent and span the affine hull of $F_y$, whose dimension
is $d$. Their barycenter
$b(y)=(d+1)^{-1}\sum_{i=0}^d p_i(y)$ lies in the relative interior
of their simplex and hence of $F_y$. Gram--Schmidt applied to
$p_i(y)-p_0(y)$ supplies a continuous definable orthonormal frame and
therefore a continuous definable family of linear isometries
\[
 A_y:K^d\longrightarrow L_y,
 \qquad L_y=\operatorname{aff}(F_y)-b(y)\subset K^n.
\]
Set $C_y=A_y^{-1}(F_y-b(y))\subset K^d$. The family $C_y$ is
Hausdorff continuous: the coordinate changes and their inverses depend
continuously on $y$, and the fibres remain bounded near each base point.
After a further shrink around $y_0$, there are constants $0<r<R$ with
\[
              r\overline B^d\subset C_y\subset R\overline B^d
                    \qquad(y\in V).                              \tag{A5.13.3}
\]
For the first inclusion use the continuously varying affine simplex
with vertices $A_y^{-1}(p_i(y)-b(y))$: its barycenter is zero, and its
finitely many face hyperplanes stay a positive distance from zero.
The second inclusion follows from Hausdorff continuity at $y_0$.

Define the radial function on $V\times S^{d-1}$ by
\[
 \rho(y,u)=\max\{t\geq0:tu\in C_y\}\in[r,R].
\]
It is definable and continuous. Here are quantitative details.
For a convex body $C$ satisfying (A5.13.3), its gauge
$p_C(v)=\inf\{a>0:v\in aC\}$ satisfies
$|p_C(v)-p_C(w)|\leq\|v-w\|/r$, by subadditivity and
$p_C(\pm(v-w))\leq\|v-w\|/r$. On the unit sphere,
$p_C\geq1/R$ and $\rho_C=1/p_C$, so $\rho_C$ is uniformly continuous.
If $C,D$ both satisfy (A5.13.3) and $d_H(C,D)\leq\delta$, then
\[
 C\subset D+\delta\overline B^d\subset(1+\delta/r)D,
 \qquad D\subset(1+\delta/r)C.
\]
The first inclusion follows from nearest points; the second uses
$r\overline B^d\subset D$ and convexity. Consequently
$\rho_C(u)\leq(1+\delta/r)\rho_D(u)$ and conversely, yielding
$|\rho_C(u)-\rho_D(u)|\leq R\delta/r$. Together with Hausdorff
continuity of $C_y$, these estimates prove joint continuity.

The typed map
\[
 H:V\times\overline B^d\longrightarrow f^{-1}(V),\qquad
 H(y,z)=
 \begin{cases}
 b(y)+A_y\bigl(\rho(y,z/\|z\|)z\bigr),&z\ne0,\\
 b(y),&z=0
 \end{cases}
\]
is a definable homeomorphism with $f\circ H=\operatorname{pr}_V$.
For fixed $y$ the formula maps each radius bijectively to the
corresponding radius of $F_y$; its inverse has coordinates
$A_y^{-1}(x-b(y))/\rho(y,u)$, with
$u=A_y^{-1}(x-b(y))/\|x-b(y)\|$ when $x\ne b(y)$, and zero at the
center. Bounds (A5.13.3) prove continuity of both maps at the center;
elsewhere it follows from the displayed formulas. A fibre at $y_0$
is definably homeomorphic to $\overline B^d$ by this same formula,
so the source can equivalently be $V\times F_{y_0}$.
This proves (A5.13.2). Its complement is nowhere dense by A5.008.
\end{proof}

% Literature check: exact searches for convex fibres + local triviality,
% definable convex fibre bundle, and convex fibres + trivialization found
% no source for this exact maximal-locus formulation. No priority claimed.
